\chapter{Produit scalaire dans le plan}

Le produit scalaire est un outil fondamental de géométrie plane qui permet de relier les longueurs, les angles et l'orthogonalité.

Il possède plusieurs expressions équivalentes (géométrique, par projection orthogonale, par les normes, analytique dans un repère orthonormé) que l'on choisit selon les données de la situation étudiée.

\section{Définition et premières expressions}

\subsection{Définition géométrique}

\begin{defin}[Norme d'un vecteur]
	\nline
	\begin{wrapstuff}[r]
		\begin{tikzpicture}[scale=1.1, >=Stealth, thick]
			\coordinate (A) at (0,0);
			\coordinate (B) at (3,1.6);
			\draw[bleu, line width=1.1pt] (A) -- (B);
			\draw[->, rouge, line width=1.3pt] (A) -- (B) node[midway, above, rouge, sloped] {$\vec{u} = \vec{AB}$};
			\fill (A) circle (1.8pt) node[below left] {$A$};
			\fill (B) circle (1.8pt) node[above right] {$B$};
		\end{tikzpicture}
	\end{wrapstuff}

	Soit $A$ et $B$ deux points du plan.

	La \textbf{norme} du vecteur $\vec{AB}$, notée $\norm{\vec{AB}}$, est la distance $AB$.

	Plus généralement, pour tout vecteur $\vec{u} = \vec{AB}$, on a $\norm{\vec{u}} = AB$.

	\rule{0mm}{12mm}
\end{defin}

\begin{defin}[Produit scalaire]
	Soit $\vec{u}$ et $\vec{v}$ deux vecteurs du plan.
	\begin{itemize}
		\item Si $\vec{u} = \vec{0}$ ou $\vec{v} = \vec{0}$, le \textbf{produit scalaire} de $\vec{u}$ par $\vec{v}$, noté $\vec{u} \cdot \vec{v}$, vaut $0$.
		\item Si $\vec{u}$ et $\vec{v}$ sont non nuls, le produit scalaire de $\vec{u}$ par $\vec{v}$ est le réel :
		      \[
			      \vec{u} \cdot \vec{v} = \norm{\vec{u}}\times\norm{\vec{v}}\times\cos\pa{\vec{u}\,;\,\vec{v}}
		      \]
	\end{itemize}
\end{defin}

\begin{rem}
	Le produit scalaire $\vec{u} \cdot \vec{v}$ se lit \og $\vec{u}$ scalaire $\vec{v}$ \fg{}.

	Il est indépendant du choix des points $A$, $B$, $C$ représentant les vecteurs, car il ne dépend que des normes et de l'angle.
\end{rem}

\begin{exemple}
	On considère deux vecteurs $\vec{u}$ et $\vec{v}$, tels que $\norm{\vec{u}} = 1$, $\norm{\vec{v}} = 2$, et $\pa{\vec{u} ; \vec{v}}=\frac{\pi}{6}$.

	On a :
	\[
		\begin{array}{rccccccl}
			\vec{u} \cdot \vec{v} & = & \norm{\vec{u}} & \times & \norm{\vec{v}} & \times & \cos\pa{\vec{u}\,;\,\vec{v}} &                                           \\
			                      & = & 1              & \times & 2              & \times & \cos(\frac{\pi}{6})          & = 2 \times \dfrac{\sqrt{3}}{2} = \sqrt{3}
		\end{array}
	\]

	\begin{center}
		\begin{tikzpicture}[scale=2, >=Stealth, thick]
			\draw[dashed, rouge, thin,domain=-1.2:1.2] plot (\x , {sqrt(3)*\x / 3});
			\draw[dashed, bleu, thin,domain=-1.2:3] plot (\x , -0.5);
			\draw[->] (0,-0.5) -- ({2},{-0.5}) node[above,sloped,midway] {$\vec{v}$};
			\draw[->] (0,0) -- ({cos(30)},{sin(30)}) node[above,sloped,midway] {$\vec{u}$};
			\coordinate (I) at ({-sqrt(3)/2}, -0.5);
			\draw[thin, black] ($(I) + (0.45,0)$) arc[start angle=0, end angle=30, radius=0.45];
			\node at ($(I) + (13:0.70)$) {\footnotesize $\frac{\pi}{6}\ \text{rad}$};
		\end{tikzpicture}
	\end{center}
\end{exemple}

\begin{attention}
	Le produit scalaire de deux vecteurs est un \textbf{nombre réel}, et non un vecteur.

	Écrire $\vec{u} \cdot \vec{v} = \vec{0}$ est une erreur : il faut écrire $\vec{u} \cdot \vec{v} = 0$.
\end{attention}

\begin{exemple}
	On considère un carré $ABCD$ de côté $3$. On souhaite calculer le produit scalaire $\vec{AB} \cdot \vec{AC}$.

	Dans le carré $ABCD$, la diagonale $[AC]$ mesure $AC = 3\sqrt{2}$ et l'angle géométrique $\widehat{BAC}$ vaut $\dfrac{\pi}{4}$ (soit $45^\circ$).

	On a alors :
	\[
		\begin{array}{rccccccl}
			\vec{AB} \cdot \vec{AC} & = & AB & \times & AC        & \times & \cos\pa{\vec{AB}\,;\,\vec{AC}} &                                            \\
			                        & = & 3  & \times & 3\sqrt{2} & \times & \cos\pa{\dfrac{\pi}{4}}        & = 9\sqrt{2} \times \dfrac{\sqrt{2}}{2} = 9
		\end{array}
	\]

	\begin{center}
		\begin{tikzpicture}[scale=1.1, >=Stealth, thick]
			\coordinate (A) at (0,0);
			\coordinate (B) at (3,0);
			\coordinate (C) at (3,3);
			\coordinate (D) at (0,3);
			\draw[black] (A) -- (B) -- (C) -- (D) -- cycle;
			\draw[thin, gris] (2.75,0) -- (2.75,0.25) -- (3,0.25);
			\draw[->, bleu, line width=1.5pt] (A) -- (B) node[midway, below, bleu] {$\vec{AB}$};
			\draw[->, rouge, line width=1.5pt] (A) -- (C) node[midway, above, sloped, rouge] {$\vec{AC}$};
			\draw[thin, black] (0.75,0) arc[start angle=0, end angle=45, radius=0.75];
			\node at (17:1.25) {$\frac{\pi}{4}\ \text{rad}$};
			\fill (A) circle (1.8pt) node[below left] {$A$};
			\fill (B) circle (1.8pt) node[below right] {$B$};
			\fill (C) circle (1.8pt) node[above right] {$C$};
			\fill (D) circle (1.8pt) node[above left] {$D$};
		\end{tikzpicture}
	\end{center}
\end{exemple}
\begin{propriete}[Carré scalaire]
	Pour tout vecteur $\vec{u} = \vec{AB}$, le produit scalaire $\vec{u} \cdot \vec{u}$ est appelé \textbf{carré scalaire} de $\vec{u}$ et noté $\vec{u}^2$.

	On a :
	\[
		\vec{u} \cdot \vec{u} = \vec{u}^2 = \norm{\vec{u}}^2 \qquad \text{et} \qquad \vec{AB} \cdot \vec{AB} = \vec{AB}^2 = AB^2
	\]
\end{propriete}

\begin{demonstration}
	Si $\vec{u} = \vec{0}$, les deux membres sont nuls.

	Sinon, l'angle entre $\vec{u}$ et lui-même est $\theta = 0$ et $\cos(0) = 1$, donc $\vec{u} \cdot \vec{u} = \norm{\vec{u}} \times \norm{\vec{u}} \times 1 = \norm{\vec{u}}^2$.
\end{demonstration}

\begin{propriete}[Signe du produit scalaire et cas particuliers]
	Soit $\vec{u}$ et $\vec{v}$ deux vecteurs non nuls formant un angle géométrique $\theta \in \intff{0}{\pi}$.

	Comme $\norm{\vec{u}}\times \norm{\vec{v}} > 0$, le produit scalaire a le signe de $\cos(\theta)$ :
	\begin{itemize}
		\item Si $\theta \in \intfo{0}{\dfrac{\pi}{2}}$ (angle aigu), alors $\vec{u} \cdot \vec{v} > 0$.

		      En particulier, si $\vec{u}$ et $\vec{v}$ sont colinéaires et de même sens ($\theta = 0$), alors $\vec{u} \cdot \vec{v} = \norm{\vec{u}}\times\norm{\vec{v}}$.
		\item Si $\theta = \dfrac{\pi}{2}$ (angle droit), alors $\vec{u} \cdot \vec{v} = 0$.
		\item Si $\theta \in \intof{\dfrac{\pi}{2}}{\pi}$ (angle obtus), alors $\vec{u} \cdot \vec{v} < 0$.

		      En particulier, si $\vec{u}$ et $\vec{v}$ sont colinéaires et de sens contraires ($\theta = \pi$), alors $\vec{u} \cdot \vec{v} = -\norm{\vec{u}}\times\norm{\vec{v}}$.
	\end{itemize}

	\begin{center}
		\begin{tikzpicture}[scale=0.85, >=Stealth, thick]
			\foreach \dx/\an/\txt in {0/50/{angle aigu \\ $\vec{u} \cdot \vec{v} > 0$}, 6/90/{angle droit \\ $\vec{u} \cdot \vec{v} = 0$}, 12/130/{angle obtus \\ $\vec{u} \cdot \vec{v} < 0$}} {
			\begin{scope}[shift={(\dx,0)}]
				\draw[->, bleu, ultra thick] (0,0) -- (2.8,0) node[below] {$\vec{u}$};
				\draw[->, rouge, ultra thick] (0,0) -- ({2.4*cos(\an)},{2.4*sin(\an)}) node[above] {$\vec{v}$};
				\draw (0.7,0) arc (0:\an:0.7);
				\node at ({1.05*cos(\an/2)},{1.05*sin(\an/2)}) {$\theta$};
				\node[align=center, below] at (1.4,-0.5) {\txt};
			\end{scope}
			}
		\end{tikzpicture}
	\end{center}
\end{propriete}

\newpage

\begin{exemple}
	On considère deux hexagones réguliers centrés en $O$ : $ABCDEF$ de côté $1$ et $A'B'C'D'E'F'$ de côté $2$, dont les sommets sont alignés avec $O$.

	\begin{center}
		\begin{tikzpicture}[scale=1.6, >=Stealth, thick]
			\coordinate (O) at (0,0);
			\foreach \i/\nom in {0/A, 60/B, 120/C, 180/D, 240/E, 300/F} { \coordinate (\nom) at (\i:1); }
			\foreach \i/\nom in {0/Ap, 60/Bp, 120/Cp, 180/Dp, 240/Ep, 300/Fp} { \coordinate (\nom) at (\i:2); }
			\foreach \i in {0, 60, 120, 180, 240, 300} { \draw[dashed, gris, thin] (O) -- (\i:2); }
			\draw[noir] (A) -- (B) -- (C) -- (D) -- (E) -- (F) -- cycle;
			\draw[noir] (Ap) -- (Bp) -- (Cp) -- (Dp) -- (Ep) -- (Fp) -- cycle;
			\draw[->, vert, line width=1.2pt] (O) -- (B) node[midway, sloped, below, vert] {$\vec{OB}$};
			\draw[->, rouge, line width=1.2pt] (O) -- (C) node[midway, sloped, above, rouge] {$\vec{OC}$};
			\draw[->, orange, line width=1.2pt] (O) -- (D) node[midway, below, orange] {$\vec{OD}$};
			\draw[->, violet, line width=1.2pt] (E) -- (Ep) node[midway, above, sloped, violet] {$\vec{EE'}$};
			\draw[->, cyan, line width=1.2pt] (F) -- (Fp) node[midway, sloped, below, cyan] {$\vec{FF'}$};
			\draw[->, teal, line width=1.2pt] (O) -- (Dp) node[near end, above, teal] {$\vec{OD'}$};
			\draw[->, magenta, line width=1.2pt] (Fp) -- (Bp) node[near start, sloped, right, below, magenta] {$\vec{F'B'}$};
			\draw[->, brown, line width=1.2pt] (Fp) -- (Ap) node[midway, below, sloped, brown] {$\vec{F'A'}$};
			\draw[->, bleu, line width=1.3pt] (O) -- (Ap) node[midway, above, bleu, fill=white] {$\vec{OA'}$};
			\node[below right, font=\footnotesize] at (A) {$A$};
			\node[above right, font=\footnotesize] at (B) {$B$};
			\node[above left, font=\footnotesize] at (C) {$C$};
			\node[below left, font=\footnotesize] at (D) {$D$};
			\node[below left, font=\footnotesize] at (E) {$E$};
			\node[below right, font=\footnotesize] at (F) {$F$};
			\node[right, font=\footnotesize] at (Ap) {$A'$};
			\node[above right, font=\footnotesize] at (Bp) {$B'$};
			\node[above left, font=\footnotesize] at (Cp) {$C'$};
			\node[left, font=\footnotesize] at (Dp) {$D'$};
			\node[below left, font=\footnotesize] at (Ep) {$E'$};
			\node[below right, font=\footnotesize] at (Fp) {$F'$};
			\fill (O) circle (1.4pt) node[below, font=\footnotesize] {$O$};
		\end{tikzpicture}
	\end{center}

	Dans un hexagone régulier, les triangles formés avec le centre $O$ sont équilatéraux, donc chaque angle au centre entre deux sommets consécutifs vaut $\dfrac{\pi}{3}$ (soit $60^\circ$). On a ainsi :

	\begin{itemize}
		\item \textbf{Angle aigu :} les vecteurs $\vec{OA'}$ et $\vec{OB}$ forment un angle de $\dfrac{\pi}{3}$, donc :
		      \[
			      \vec{OA'} \cdot \vec{OB} = OA' \times OB \times \cos\pa{\dfrac{\pi}{3}} = 2 \times 1 \times \dfrac{1}{2} = 1 > 0
		      \]

		\item \textbf{Angle obtus :} les vecteurs $\vec{OA'}$ et $\vec{OC}$ forment un angle de $\dfrac{2\pi}{3}$, donc :
		      \[
			      \vec{OA'} \cdot \vec{OC} = OA' \times OC \times \cos\pa{\dfrac{2\pi}{3}} = 2 \times 1 \times \pa{-\dfrac{1}{2}} = -1 < 0
		      \]

		\item \textbf{Vecteurs colinéaires de sens contraires :} les vecteurs $\vec{OA'}$ et $\vec{OD}$ forment un angle plat $\pi$, donc :
		      \[
			      \vec{OA'} \cdot \vec{OD} = -OA' \times OD = -2 \times 1 = -2 < 0
		      \]
		\item \textbf{Vecteurs colinéaires de même sens :} les vecteurs $\vec{EE'}$ et $\vec{FF'}$ ont des directions faisant un angle de $\dfrac{\pi}{3}$, avec $EE' = 1$ et $FF' = 1$, donc :
		      \[
			      \vec{EE'} \cdot \vec{FF'} = EE' \times FF' \times \cos\pa{\dfrac{\pi}{3}} = 1 \times 1 \times \dfrac{1}{2} = \dfrac{1}{2}
		      \]
		\item \textbf{Vecteurs orthogonaux :} En revanche, la droite $(OD')$ est perpendiculaire à $(F'B'$, donc $\vec{OD'}$ et $\vec{F'B'}$ sont orthogonaux :
		      \[
			      \vec{OD'} \cdot \vec{F'B'} = 0
		      \]
		\item \textbf{Vecteurs n'ayant pas la même origine :} le triangle $OF'A'$ étant équilatéral de côté $2$, le vecteur $\vec{F'A'}$ est égal au vecteur $\vec{OB'}$ (de norme $2$). Les directions de $\vec{OD'}$ et $\vec{F'A'}$ forment donc un angle de $\dfrac{2\pi}{3}$, d'où :
		      \[
			      \vec{OD'} \cdot \vec{F'A'} = OD' \times F'A' \times \cos\pa{\dfrac{2\pi}{3}} = 2 \times 2 \times \pa{-\dfrac{1}{2}} = -2
		      \]
	\end{itemize}
\end{exemple}

\newpage

\begin{methode}[Calculer un produit scalaire avec la formule du cosinus]
	Pour calculer $\vec{u} \cdot \vec{v}$ à l'aide de la définition géométrique :
	\begin{enumerate}
		\item on détermine les normes $\norm{\vec{u}}$ et $\norm{\vec{v}}$ ;
		\item on ramène les deux vecteurs à la \textbf{même origine} afin d'identifier l'angle géométrique $\theta$ formé par leurs directions ;
		\item on applique $\vec{u} \cdot \vec{v} = \norm{\vec{u}}\times\norm{\vec{v}}\cos(\theta)$ et on vérifie la cohérence du signe.
	\end{enumerate}

	\ms
	Soit $ABC$ un triangle équilatéral de côté $5$ et $I$ le milieu de $[AB]$.

	Calculer $\vec{AB} \cdot \vec{AC}$ et $\vec{AI} \cdot \vec{BC}$.

	\begin{correction}
		\nline
		\begin{enumerate}
			\item Les vecteurs $\vec{AB}$ et $\vec{AC}$ ont la même origine $A$, avec $AB = AC = 5$ et $\widehat{BAC} = \dfrac{\pi}{3}$ :
			      \[
				      \vec{AB} \cdot \vec{AC} = 5 \times 5 \times \cos\pa{\dfrac{\pi}{3}} = 25 \times \dfrac{1}{2} = \dfrac{25}{2} = 12{,}5
			      \]
			\item Les vecteurs $\vec{AI}$ et $\vec{BC}$ n'ont pas la même origine.

			      On construit le point $D$ tel que $\vec{AD} = \vec{BC}$.

			      Alors $\widehat{IAD} = \pi - \dfrac{\pi}{3} = \dfrac{2\pi}{3}$, $AI = 2{,}5$ et $AD = BC = 5$ :
			      \[
				      \vec{AI} \cdot \vec{BC} = \vec{AI} \cdot \vec{AD} = 2{,}5 \times 5 \times \cos\pa{\dfrac{2\pi}{3}} = 12{,}5 \times \pa{-\dfrac{1}{2}} = -6{,}25
			      \]
		\end{enumerate}

		\begin{center}
			\begin{tikzpicture}[scale=0.65, >=Stealth, thick]
				\coordinate (A) at (0,0);
				\coordinate (B) at (5,0);
				\coordinate (C) at (60:5);
				\coordinate (I) at (2.5,0);
				\coordinate (D) at (120:5);
				\draw[->, rouge, line width=1.5pt] (A) -- (D) node[midway, below, sloped] {$\vec{AD} = \vec{BC}$};
				\draw[vert, fill=vert!15] (0.9,0) arc[start angle=0, end angle=120, radius=0.9] -- (A) -- cycle;
				\node[vert, above] at (90:0.9) {$\frac{2\pi}{3}$};
				\node[below left] at (A) {$A$};
				\node[below right] at (B) {$B$};
				\node[above] at (C) {$C$};
				\node[below] at (I) {$I$};
				\node[above left] at (D) {$D$};
				\draw[black] (A) -- (B) -- (C) -- cycle;
				\draw[->, rouge, line width=1.5pt] (A) -- (I) node[midway, below] {$\vec{AI}$};
				\draw[->, rouge, line width=1.5pt] (B) -- (C) node[midway, above, sloped] {$\vec{BC}$};
				\foreach \pt in {A,B,C,D,I} \filldraw (\pt) circle (2pt);
			\end{tikzpicture}
		\end{center}
	\end{correction}
\end{methode}

\subsection{Expression par projection orthogonale}

\begin{defin}[Projeté orthogonal sur une droite]
	\nline
	\begin{wrapstuff}[r]
		\begin{tikzpicture}[x=7mm, y=7mm, >=Stealth, thick]
			\coordinate (D1) at (-1.5, -0.45);
			\coordinate (D2) at (4.5, 1.35);
			\draw[bleu, line width=1.1pt] (D1) -- (D2) node[above left, bleu] {$(d)$};
			\coordinate (H) at (1.5, 0.45);
			\coordinate (M) at ($(H) + (-0.6, 2.0)$);
			\draw[dashed, gris] ($(H) - (-0.21, 0.7)$) -- ($(M) + (-0.21, 0.7)$);
			\coordinate (U) at ($(H) + (0.28, 0.084)$);
			\coordinate (V) at ($(H) + (-0.084, 0.28)$);
			\draw[thin, noir] (U) -- ($(U) + (-0.084, 0.28)$) -- (V);
			\fill[rouge] (M) circle (1.8pt) node[above right, rouge] {$M$};
			\fill[vert] (H) circle (1.8pt) node[below right, vert] {$H$};
		\end{tikzpicture}
	\end{wrapstuff}

	Soit $(d)$ une droite et $M$ un point du plan.

	Le \textbf{projeté orthogonal} de $M$ sur $(d)$ est le point d'intersection $H$ de la droite $(d)$ et de la perpendiculaire à $(d)$ passant par $M$.

	\rule{0mm}{10mm}
\end{defin}

\begin{propriete}[Produit scalaire et projection orthogonale]
	\nline
	\begin{wrapstuff}[r]
		\begin{tikzpicture}[x=7mm, y=7mm, >=Stealth, thick]
			\coordinate (O) at (0, 0);
			\coordinate (A) at (4.5, 1.35);
			\coordinate (H) at (2.8, 0.84);
			\coordinate (B) at ($(H) + (-0.66, 2.2)$);
			\draw[dashed, gris, thin] (-1.2, -0.36) -- (5.5, 1.65);
			\draw[dashed, gris] (B) -- (H);
			\draw[thin, noir] ($(H) + (0.28, 0.084)$) -- ($(H) + (0.196, 0.364)$) -- ($(H) + (-0.084, 0.28)$);
			\draw[->, bleu, line width=1.2pt] (O) -- (A);
			\draw[->, rouge, line width=1.2pt] (O) -- (B);
			\draw[->, vert, line width=1.2pt] (O) -- (H);
			\fill[noir] (O) circle (1.8pt) node[below, noir] {$O$};
			\fill[bleu] (A) circle (1.8pt) node[below, bleu] {$A$};
			\fill[rouge] (B) circle (1.8pt) node[above, rouge] {$B$};
			\fill[vert] (H) circle (1.8pt) node[below, vert] {$H$};
		\end{tikzpicture}
	\end{wrapstuff}

	Soit $O$, $A$ et $B$ trois points du plan tels que $O \neq A$, et soit $H$ le projeté orthogonal de $B$ sur la droite $(OA)$.

	Alors :
	\[
		\vec{OA} \cdot \vec{OB} = \vec{OA} \cdot \vec{OH} =
		\begin{cases}
			OA \times OH  & \,\text{si }\,\vec{OA} \text{ et } \vec{OH}\,\text{ sont de même sens,}       \\
			-OA \times OH & \,\text{si }\,\vec{OA} \text{ et } \vec{OH}\,\text{ sont de sens contraires.}
		\end{cases}
	\]
\end{propriete}

\newpage

\begin{demonstration}
	Si $B$ est sur $(OA)$, alors $H = B$ et le résultat est le cas de vecteurs colinéaires.

	Si $\vec{OB} = \vec{0}$ ou si $H = O$, alors $\vec{OA}$ et $\vec{OB}$ sont orthogonaux, les deux produits sont nuls.

	Sinon, soit $\theta = \widehat{AOB}$ :
	\begin{itemize}
		\item Si $\theta$ est aigu, $H$ est sur la demi-droite $[OA)$.

		      Dans le triangle $OHB$ rectangle en $H$, on a : $OH = OB\cos(\theta)$.

		      Alors :
		      \[
			      \vec{OA} \cdot \vec{OB} = OA \times OB \times \cos(\theta) = OA \times OH
		      \]

		      \ldots avec $\vec{OA}$, $\vec{OH}$ de même sens.
		\item Si $\theta$ est obtus, $H$ est de l'autre côté de $O$.

		      On a $OH = OB\cos(\pi - \theta) = -OB\cos(\theta)$.

		      Alors :

		      \[
			      \vec{OA} \cdot \vec{OB} = OA \times OB \times \cos(\theta) = -OA \times OH
		      \]

		      \ldots avec $\vec{OA}$, $\vec{OH}$ de sens contraires.
	\end{itemize}

	Dans les deux cas, $\vec{OA} \cdot \vec{OH}$ vaut respectivement $OA \times OH$ et $-OA \times OH$, d'où l'égalité $\vec{OA} \cdot \vec{OB} = \vec{OA} \cdot \vec{OH}$.

	\begin{center}
		\begin{tikzpicture}[x=9mm, y=9mm, >=Stealth, thick]
			\begin{scope}[shift={(0,0)}]
				\coordinate (O) at (0,0);
				\coordinate (A) at (4.5,0);
				\coordinate (B) at (3,2.2);
				\coordinate (H) at (3,0);
				\draw[->, bleu, line width=1.2pt] (O) -- (A) node[right] {$A$};
				\draw[->, rouge, line width=1.2pt] (O) -- (B) node[above right] {$B$};
				\draw[dashed, gris] (B) -- (H);
				\draw (H) ++(0,0.25) -- ++(0.25,0) -- ++(0,-0.25);
				\draw[vert, ultra thick] (0,-0.16) -- (3,-0.16);
				\draw[thin, noir] (0.8,0) arc[start angle=0, end angle=36.25, radius=0.8];
				\node[font=\footnotesize] at (18.1:1.1) {$\theta$};
				\fill[rouge] (B) circle (2pt);
				\fill (O) circle (2pt) node[below left] {$O$};
				\fill[vert] (H) circle (2pt) node[below right, vert] {$H$};
				\node[font=\footnotesize, below] at (2.25,-0.6) {$\vec{OA} \cdot \vec{OB} = OA \times OH > 0$};
			\end{scope}
			\begin{scope}[shift={(7.5,0)}]
				\coordinate (O) at (1.5,0);
				\coordinate (A) at (5,0);
				\coordinate (B) at (0,2.2);
				\coordinate (H) at (0,0);
				\draw[dashed, gris] (-0.4,0) -- (O);
				\draw[->, bleu, line width=1.2pt] (O) -- (A) node[right] {$A$};
				\draw[->, rouge, line width=1.2pt] (O) -- (B) node[above left] {$B$};
				\draw[dashed, gris] (B) -- (H);
				\draw (H) ++(0,0.25) -- ++(0.25,0) -- ++(0,-0.25);
				\draw[vert, ultra thick] (1.5,-0.16) -- (0,-0.16);
				\fill (O) circle (2pt) node[below right] {$O$};
				\fill[rouge] (B) circle (2pt);
				\draw[thin, noir] ($(O) + (0.65,0)$) arc[start angle=0, end angle=124.29, radius=0.65];
				\node[font=\footnotesize] at ($(O) + (62.1:0.92)$) {$\theta$};
				\fill[vert] (H) circle (2pt) node[below left, vert] {$H$};
				\node[font=\footnotesize, below] at (2.5,-0.6) {$\vec{OA} \cdot \vec{OB} = -OA \times OH < 0$};
			\end{scope}
		\end{tikzpicture}
	\end{center}
\end{demonstration}

\begin{methode}[Calculer un produit scalaire par projection orthogonale]
	Soit $ABCD$ un carré de côté $4$.

	Calculer $\vec{AB} \cdot \vec{AC}$, $\vec{AB} \cdot \vec{AD}$ et $\vec{AD} \cdot \vec{CB}$.

	\begin{correction}
		\nline
		\begin{wrapstuff}[r]
			\begin{tikzpicture}[x=8mm, y=8mm, >=Stealth, thick]
				\coordinate (A) at (0,0);
				\coordinate (B) at (4,0);
				\coordinate (C) at (4,4);
				\coordinate (D) at (0,4);
				\draw[noir] (A) -- (B) -- (C) -- (D) -- cycle;
				\draw[thin, gris] (3.7,0) -- (3.7,0.3) -- (4,0.3);
				\draw[->, bleu, line width=1.5pt] (A) -- (B);
				\draw[->, rouge, line width=1.5pt] (A) -- (C);
				\draw[->, vert, line width=1.5pt] (A) -- (D);
				\draw[->, orange, line width=1.5pt] (C) -- (B);
				\node[below, font=\footnotesize] at (2,-0.15) {$4$};
				\fill (A) circle (1.8pt) node[below left] {$A$};
				\fill (B) circle (1.8pt) node[below right] {$B$};
				\fill (C) circle (1.8pt) node[above right] {$C$};
				\fill (D) circle (1.8pt) node[above left] {$D$};
			\end{tikzpicture}
		\end{wrapstuff}
		\begin{enumerate}
			\item Le point $B$ est le projeté orthogonal de $C$ sur la droite $(AB)$, donc :
			      \[
				      \def\arraystretch{1.5}
				      \begin{array}{rcl}
					      \vec{AB} \cdot \vec{AC} & = & \vec{AB} \cdot \vec{AB} \\
					                              & = & AB^2 = 4^2 = 16
				      \end{array}
			      \]
			\item Les côtés $[AB]$ et $[AD]$ sont perpendiculaires, donc $\vec{AB} \cdot \vec{AD} = 0$.
			\item Comme $ABCD$ est un carré, $\vec{CB} = \vec{DA} = -\vec{AD}$, d'où :
			      \[
				      \def\arraystretch{1.5}
				      \begin{array}{rcl}
					      \vec{AD} \cdot \vec{CB} & = & \vec{AD} \cdot \pa{-\vec{AD}} \\
					                              & = & -AD^2                         \\
					                              & = & -4^2 = -16
				      \end{array}
			      \]
		\end{enumerate}
	\end{correction}
\end{methode}

\newpage

\section{Propriétés algébriques et formules de polarisation}

\subsection{Symétrie et bilinéarité}

\begin{propriete}[Symétrie et bilinéarité du produit scalaire]
	Pour tous vecteurs $\vec{u}$, $\vec{v}$, $\vec{w}$ du plan et tout nombre réel $k$ :
	\begin{itemize}
		\item \textbf{Symétrie :}
		      \begin{itemize}[itemsep=2mm]
			      \item $\vec{u} \cdot \vec{v} = \vec{v} \cdot \vec{u}$
		      \end{itemize}
		\item \textbf{Bilinéarité :}
		      \begin{itemize}[itemsep=2mm]
			      \item $\vec{u} \cdot \pa{\vec{v} + \vec{w}} = \vec{u} \cdot \vec{v} + \vec{u} \cdot \vec{w}$
			      \item $\pa{\vec{u} + \vec{v}} \cdot \vec{w} = \vec{u} \cdot \vec{w} + \vec{v} \cdot \vec{w}$
			      \item $\pa{k\vec{u}} \cdot \vec{v} = \vec{u} \cdot \pa{k\vec{v}} = k\pa{\vec{u} \cdot \vec{v}}$
		      \end{itemize}
	\end{itemize}
\end{propriete}

\begin{demonstration}
	\nline
	\begin{itemize}
		\item \textbf{Symétrie :}

		      Si $\vec{u} = \vec{0}$ ou $\vec{v} = \vec{0}$, alors $\vec{u} \cdot \vec{v} = 0 = \vec{v} \cdot \vec{u}$.

		      Si $\vec{u} \neq \vec{0}$ et $\vec{v} \neq \vec{0}$, les angles géométriques $\pa{\vec{u}\,;\,\vec{v}}$ et $\pa{\vec{v}\,;\,\vec{u}}$ sont égaux, donc par définition :
		      \[
			      \vec{u} \cdot \vec{v} = \norm{\vec{u}}\norm{\vec{v}}\cos\pa{\vec{u}\,;\,\vec{v}} = \norm{\vec{v}}\norm{\vec{u}}\cos\pa{\vec{v}\,;\,\vec{u}} = \vec{v} \cdot \vec{u}
		      \]

		\item \textbf{Multiplication par un réel $k$ :}

		      Si $k = 0$, $\vec{u} = \vec{0}$ ou $\vec{v} = \vec{0}$, les trois membres valent $0$.

		      Supposons désormais $k \neq 0$, $\vec{u} \neq \vec{0}$ et $\vec{v} \neq \vec{0}$ :
		      \begin{itemize}
			      \item Si $k > 0$, alors $\norm{k\vec{u}} = k\norm{\vec{u}}$ et $k\vec{u}$ a le même sens que $\vec{u}$, donc $\cos\pa{k\vec{u}\,;\,\vec{v}} = \cos\pa{\vec{u}\,;\,\vec{v}}$. D'où :
			            \[
				            \pa{k\vec{u}} \cdot \vec{v} = \norm{k\vec{u}}\times \norm{\vec{v}}\cos\pa{k\vec{u}\,;\,\vec{v}} = k\norm{\vec{u}}\times\norm{\vec{v}}\cos\pa{\vec{u}\,;\,\vec{v}} = k\pa{\vec{u} \cdot \vec{v}}
			            \]
			      \item Si $k < 0$, alors $\norm{k\vec{u}} = -k\norm{\vec{u}}$ et $k\vec{u}$ est de sens contraire à $\vec{u}$, donc $\pa{k\vec{u}\,;\,\vec{v}} = \pi - \pa{\vec{u}\,;\,\vec{v}}$, ce qui donne $\cos\pa{k\vec{u}\,;\,\vec{v}} = -\cos\pa{\vec{u}\,;\,\vec{v}}$. D'où :
			            \[
				            \pa{k\vec{u}} \cdot \vec{v} = \pa{-k\norm{\vec{u}}}\times\norm{\vec{v}}\pa{-\cos\pa{\vec{u}\,;\,\vec{v}}} = k\norm{\vec{u}}\times\norm{\vec{v}}\cos\pa{\vec{u}\,;\,\vec{v}} = k\pa{\vec{u} \cdot \vec{v}}
			            \]
		      \end{itemize}

		      Par symétrie, on obtient également $\vec{u} \cdot \pa{k\vec{v}} = \pa{k\vec{v}} \cdot \vec{u} = k\pa{\vec{v} \cdot \vec{u}} = k\pa{\vec{u} \cdot \vec{v}}$.
	\end{itemize}

	\begin{itemize}
		\item \textbf{Distributivité :}

		      Si $\vec{u} = \vec{0}$, l'égalité est évidente.

		      Supposons $\vec{u} \neq \vec{0}$ et posons $\vec{u} = \vec{AB}$, $\vec{v} = \vec{BC}$ et $\vec{w} = \vec{CD}$.
		      \begin{center}
			      \begin{tikzpicture}[>=Stealth, thick]
				      \coordinate (A) at (1,0);
				      \coordinate (B) at (3.2,0);
				      \coordinate (C) at (4.7,1.8);
				      \coordinate (D) at (6.2,1.1);
				      \coordinate (Cp) at (4.7,0);
				      \coordinate (Dp) at (6.2,0);
				      \draw[dashed, gris] (0.2,0) -- (7.2,0);
				      \draw[dashed, gris] (C) -- (Cp);
				      \draw[dashed, gris] (D) -- (Dp);
				      \draw[thin, noir] (4.75,0.25) -- (5,0.25) -- (5,0);
				      \draw[thin, noir] (6.2,0.25) -- (6.45,0.25) -- (6.45,0);
				      \draw[->, bleu, line width=1.2pt] (A) -- (B) node[sloped, midway, below, bleu] {$\vec{u}$};
				      \draw[->, rouge, line width=1.2pt] (B) -- (C) node[sloped, midway, above, rouge] {$\vec{v}$};
				      \draw[->, vert, line width=1.2pt] (C) -- (D) node[sloped, midway, above, vert] {$\vec{w}$};
				      \draw[->, orange, line width=1.1pt] (B) -- (D) node[sloped, midway, above, orange, fill=white] {$\vec{v}+\vec{w}$};
				      \draw[->, rouge, line width=1.1pt] (B) -- (Cp);% node[sloped, midway, below, rouge] {\footnotesize $\vec{BC'}$};
				      \draw[->, vert, line width=1.1pt] (Cp) -- (Dp);% node[sloped, midway, below, vert] {\footnotesize $\vec{C'D'}$};
				      \fill (A) circle (1.6pt) node[below] {$A$};
				      \fill (B) circle (1.6pt) node[below] {$B$};
				      \fill (C) circle (1.6pt) node[above] {$C$};
				      \fill (D) circle (1.6pt) node[above] {$D$};
				      \fill[rouge] (Cp) circle (1.6pt) node[below, rouge] {$C'$};
				      \fill[vert] (Dp) circle (1.6pt) node[below, vert] {$D'$};
			      \end{tikzpicture}
		      \end{center}
		      D'après la relation de Chasles, on a $\vec{v} + \vec{w} = \vec{BC} + \vec{CD} = \vec{BD}$.

		      Notons $C'$ et $D'$ les projetés orthogonaux respectifs de $C$ et $D$ sur la droite $(AB)$.

		      Par propriété du produit scalaire et de la projection orthogonale :
		      \begin{itemize}
			      \item $\vec{u} \cdot \vec{v} = \vec{AB} \cdot \vec{BC} = \vec{AB} \cdot \vec{BC'}$
			      \item $\vec{u} \cdot \vec{w} = \vec{AB} \cdot \vec{CD} = \vec{AB} \cdot \vec{C'D'}$
			      \item $\vec{u} \cdot \pa{\vec{v} + \vec{w}} = \vec{AB} \cdot \vec{BD} = \vec{AB} \cdot \vec{BD'}$
		      \end{itemize}

		      Comme $C'$ et $D'$ appartiennent à la droite $(AB)$, les vecteurs $\vec{AB}$, $\vec{BC'}$ et $\vec{C'D'}$ sont colinéaires à $(AB)$, et par relation de Chasles on a $\vec{BC'} + \vec{C'D'} = \vec{BD'}$.

		      En étudiant les sens de ces vecteurs colinéaires sur la droite $(AB)$ (avec $\vec{AB}$ non nul), on obtient :
		      \[
			      \vec{AB} \cdot \vec{BC'} + \vec{AB} \cdot \vec{C'D'} = \vec{AB} \cdot \vec{BD'}
		      \]

		      Soit :
		      \[
			      \vec{u} \cdot \vec{v} + \vec{u} \cdot \vec{w} = \vec{u} \cdot \pa{\vec{v} + \vec{w}}
		      \]

		      Enfin, par symétrie du produit scalaire :
		      \[
			      \pa{\vec{u} + \vec{v}} \cdot \vec{w} = \vec{w} \cdot \pa{\vec{u} + \vec{v}} = \vec{w} \cdot \vec{u} + \vec{w} \cdot \vec{v} = \vec{u} \cdot \vec{w} + \vec{v} \cdot \vec{w}
		      \]
	\end{itemize}
\end{demonstration}

\begin{propriete}[Identités remarquables scalaires]
	Pour tous vecteurs $\vec{u}$ et $\vec{v}$ du plan :
	\begin{enumerate}
		\item $\norm{\vec{u} + \vec{v}}^2 = \pa{\vec{u} + \vec{v}}^2 = \norm{\vec{u}}^2 + 2\,\vec{u} \cdot \vec{v} + \norm{\vec{v}}^2$ ;
		\item $\norm{\vec{u} - \vec{v}}^2 = \pa{\vec{u} - \vec{v}}^2 = \norm{\vec{u}}^2 - 2\,\vec{u} \cdot \vec{v} + \norm{\vec{v}}^2$ ;
		\item $\pa{\vec{u} + \vec{v}} \cdot \pa{\vec{u} - \vec{v}} = \norm{\vec{u}}^2 - \norm{\vec{v}}^2$.
	\end{enumerate}
\end{propriete}

\begin{demonstration}
	\nline
	\begin{wrapstuff}[r]
		\begin{tikzpicture}[x=8mm, y=8mm, >=Stealth, thick]
			\coordinate (O) at (0,0);
			\coordinate (A) at (3,0.5);
			\coordinate (B) at (1,3);
			\coordinate (C) at (4,3.5);
			\draw[dashed, gris] (A) -- (C) -- (B);
			\draw[->, bleu, ultra thick] (O) -- (A) node[sloped, midway, below, font=\footnotesize] {$\vec{u}$};
			\draw[->, rouge, ultra thick] (O) -- (B) node[sloped, midway, above, font=\footnotesize] {$\vec{v}$};
			\draw[->, vert, ultra thick] (O) -- (C) node[sloped, pos=0.7, below, font=\footnotesize] {$\vec{u} + \vec{v}$};
			\draw[->, orange, ultra thick] (B) -- (A) node[sloped, pos=0.52, below left, font=\footnotesize] {$\vec{u} - \vec{v}$};
			\fill (O) circle (2pt) node[below left] {$O$};
		\end{tikzpicture}
	\end{wrapstuff}

	Par bilinéarité et symétrie du produit scalaire :

	\begin{enumerate}
		\item $\norm{\vec{u} + \vec{v}}^2 = \vec{u} \cdot \vec{u} + \vec{u} \cdot \vec{v} + \vec{v} \cdot \vec{u} + \vec{v} \cdot \vec{v} = \norm{\vec{u}}^2 + 2\,\vec{u} \cdot \vec{v} + \norm{\vec{v}}^2$
		\item $\norm{\vec{u} - \vec{v}}^2 = \vec{u} \cdot \vec{u} - \vec{u} \cdot \vec{v} - \vec{v} \cdot \vec{u} + \vec{v} \cdot \vec{v} = \norm{\vec{u}}^2 - 2\,\vec{u} \cdot \vec{v} + \norm{\vec{v}}^2$
		\item $\pa{\vec{u} + \vec{v}} \cdot \pa{\vec{u} - \vec{v}}= \vec{u} \cdot \vec{u} - \vec{u} \cdot \vec{v} + \vec{v} \cdot \vec{u} - \vec{v} \cdot \vec{v} = \norm{\vec{u}}^2 - \norm{\vec{v}}^2$
	\end{enumerate}

	\rule{0mm}{10mm}
\end{demonstration}

\begin{methode}[Développer et calculer avec les propriétés algébriques]
	Soit $\vec{u}$ et $\vec{v}$ deux vecteurs tels que $\norm{\vec{u}} = 2$, $\norm{\vec{v}} = 3$ et $\vec{u} \cdot \vec{v} = 1$.

	Calculer :
	\begin{enumerate}
		\item $\pa{\vec{u} + \vec{v}} \cdot \pa{\vec{u} - \vec{v}}$ ;
		\item $\norm{\vec{u} + \vec{v}}$ ;
		\item $\pa{-2\vec{v}} \cdot \pa{3\vec{u} - \vec{v}}$.
	\end{enumerate}

	\begin{correction}
		\nline
		\begin{enumerate}
			\item $\pa{\vec{u} + \vec{v}} \cdot \pa{\vec{u} - \vec{v}} = \norm{\vec{u}}^2 - \norm{\vec{v}}^2 = 2^2 - 3^2 = -5$.
			\item $\norm{\vec{u} + \vec{v}}^2 = \norm{\vec{u}}^2 + 2\,\vec{u} \cdot \vec{v} + \norm{\vec{v}}^2 = 4 + 2 + 9 = 15$, donc $\norm{\vec{u} + \vec{v}} = \sqrt{15}$.
			\item $\pa{-2\vec{v}} \cdot \pa{3\vec{u} - \vec{v}} = -6\pa{\vec{v} \cdot \vec{u}} + 2\norm{\vec{v}}^2 = -6 \times 1 + 2 \times 9 = 12$.
		\end{enumerate}
	\end{correction}
\end{methode}

\subsection{Formules de polarisation (expression avec les normes)}

\begin{propriete}[Formules de polarisation]
	Pour tous vecteurs $\vec{u}$ et $\vec{v}$ du plan :
	\vspace{-3mm}
	\begin{multicols}{2}
		\begin{itemize}
			\item $\vec{u} \cdot \vec{v} = \dfrac{1}{2}\pa{\norm{\vec{u}}^2 + \norm{\vec{v}}^2 - \norm{\vec{u} - \vec{v}}^2}$
			\item $\vec{u} \cdot \vec{v} = \dfrac{1}{2}\pa{\norm{\vec{u} + \vec{v}}^2 - \norm{\vec{u}}^2 - \norm{\vec{v}}^2}$
		\end{itemize}
	\end{multicols}

	En particulier, pour tous points $A$, $B$ et $C$ (en posant $\vec{u} = \vec{AB}$ et $\vec{v} = \vec{AC}$, d'où $\vec{u} - \vec{v} = \vec{CB}$) :
	\[
		\vec{AB} \cdot \vec{AC} = \dfrac{1}{2}\pa{AB^2 + AC^2 - BC^2}
	\]
\end{propriete}

\begin{demonstration}
	D'après les identités remarquables scalaires :
	\[
		\def\arraystretch{1.5}
		\begin{array}{rcl}
			\norm{\vec{u} + \vec{v}}^2 = \norm{\vec{u}}^2 + 2\,\vec{u} \cdot \vec{v} + \norm{\vec{v}}^2 & \eqv & 2\,\vec{u} \cdot \vec{v} = \norm{\vec{u} + \vec{v}}^2 - \norm{\vec{u}}^2 - \norm{\vec{v}}^2 \\
			\norm{\vec{u} - \vec{v}}^2 = \norm{\vec{u}}^2 - 2\,\vec{u} \cdot \vec{v} + \norm{\vec{v}}^2 & \eqv & 2\,\vec{u} \cdot \vec{v} = \norm{\vec{u}}^2 + \norm{\vec{v}}^2 - \norm{\vec{u} - \vec{v}}^2
		\end{array}
	\]
	En divisant par $2$, on obtient les deux expressions.
\end{demonstration}

\begin{exemple}
	\nline
	\begin{wrapstuff}[r]
		\begin{tikzpicture}[x=5mm, y=5mm, >=Stealth, thick]
			\coordinate (C) at (0,0);
			\coordinate (F) at (7,0);
			\coordinate (G) at ({6*19/21},{6*4*sqrt(5)/21});

			\draw[noir] (C) -- (F) -- (G) -- cycle;
			\draw[->, bleu, line width=1.2pt] (C) -- (G) node[midway, above, sloped, bleu] {$\vec{CG}$};
			\draw[->, rouge, line width=1.2pt] (C) -- (F) node[midway, below, sloped, rouge] {$\vec{CF}$};
			\fill (C) circle (2pt) node[below left] {$C$};
			\fill (F) circle (2pt) node[below right] {$F$};
			\fill (G) circle (2pt) node[above] {$G$};
		\end{tikzpicture}
	\end{wrapstuff}

	Soit $CFG$ un triangle tel que $CG = 6$, $CF = 7$ et $GF = 3$.

	Alors :
	\[
		\vec{CG} \cdot \vec{CF} = \dfrac{1}{2}\pa{CG^2 + CF^2 - GF^2} = \dfrac{1}{2}\pa{36 + 49 - 9} = \dfrac{76}{2} = 38
	\]
\end{exemple}

\section{Produit scalaire dans un repère orthonormé}

Dans toute cette section, le plan est muni d'un repère orthonormé $\vOIJ$.

\begin{propriete}[Expression analytique du produit scalaire]
	Soit $\vec{u}\coord{x}{y}$ et $\vec{v}\coord{x'}{y'}$ deux vecteurs du plan.

	Alors :
	\[
		\vec{u} \cdot \vec{v} = xx' + yy' \qquad \text{et} \qquad \norm{\vec{u}} = \sqrt{x^2 + y^2}
	\]

	\begin{center}
		\begin{tikzpicture}[x=11mm, y=11mm, >=Stealth, thick]
			\draw[->, gris] (-0.4,0) -- (3.8,0);
			\draw[->, gris] (0,-0.4) -- (0,3.2);
			\draw[->, noir, line width=1.1pt] (0,0) -- (1,0) node[midway, sloped, below] {$\vec{\imath}$};
			\draw[->, noir, line width=1.1pt] (0,0) -- (0,1) node[midway, sloped, above] {$\vec{\jmath}$};
			\coordinate (U) at (3,1.5);
			\coordinate (V) at (1.25,2.5);
			\draw[dashed, bleu, thin] (3,0) -- (U) -- (0,1.5);
			\draw[dashed, rouge, thin] (1.25,0) -- (V) -- (0,2.5);
			\draw[->, bleu, line width=1.3pt] (0,0) -- (U) node[pos=0.55, sloped, above, bleu] {$\vec{u}$};
			\draw[->, rouge, line width=1.3pt] (0,0) -- (V) node[midway, sloped, below, rouge] {$\vec{v}$};
			\node[below, bleu] at (3,0) {$x$};
			\node[left, bleu] at (0,1.5) {$y$};
			\node[below, rouge] at (1.25,0) {$x'$};
			\node[left, rouge] at (0,2.5) {$y'$};
			\fill (0,0) circle (1.5pt) node[below left] {$O$};
		\end{tikzpicture}
	\end{center}
\end{propriete}

\begin{demonstration}
	Le repère étant orthonormé, $\norm{\vec{i}} = \norm{\vec{j}} = 1$ et $\vec{i} \cdot \vec{j} = 0$ (les vecteurs sont orthogonaux).

	En décomposant $\vec{u} = x\vec{i} + y\vec{j}$ et $\vec{v} = x'\vec{i} + y'\vec{j}$, la bilinéarité donne :
	\[
		\vec{u} \cdot \vec{v} = \pa{x\vec{i} + y\vec{j}}\cdot\pa{x'\vec{i} + y'\vec{j}}
	\]

	Donc :
	\[
		\begin{array}{rccccccccl}
			\vec{u} \cdot \vec{v} & = & xx'\,\vec{i}^2 & + & xy'\pa{\vec{i} \cdot \vec{j}} & + & yx'\pa{\vec{j} \cdot \vec{i}} & + & yy'\,\vec{j}^2 &             \\
			                      & = & xx'            & + & 0                             & + & 0                             & + & yy'            & = xx' + yy'
		\end{array}
	\]

	En particulier, $\norm{\vec{u}}^2 = \vec{u} \cdot \vec{u} = x^2 + y^2$, d'où $\norm{\vec{u}} = \sqrt{x^2 + y^2}$.
\end{demonstration}

\begin{attention}
	Cette formule n'est valable \textbf{que dans un repère orthonormé}.

	Si le repère est quelconque, on peut avoir $\vec{i} \cdot \vec{j} \neq 0$ ou des vecteurs de base de norme différente de $1$, et la formule est fausse.
\end{attention}

\newpage

\begin{propriete}[Coordonnées et distance]
	Pour tout vecteur $\vec{u}\coord{x}{y}$ :
	\[
		x = \vec{u} \cdot \vec{i}\qquad\text{et}\qquad y = \vec{u} \cdot \vec{j}
	\]

	De plus, pour $A\coordl{x_A}{y_A}$ et $B\coordl{x_B}{y_B}$ :
	\[
		AB = \sqrt{\pa{x_B - x_A}^2 + \pa{y_B - y_A}^2}
	\]
\end{propriete}

\begin{demonstration}
	Par l'expression analytique :
	\[
		\vec{u} \cdot \vec{i} = x \times 1 + y \times 0 = x\qquad\text{et}\qquad\vec{u} \cdot \vec{j} = x \times 0 + y \times 1 = y
	\]

	De plus, $\vec{AB}\coord{x_B - x_A}{y_B - y_A}$ et $AB = \norm{\vec{AB}}=\sqrt{\pa{x_B - x_A}^2 + \pa{y_B - y_A}^2}$.
\end{demonstration}

\begin{methode}[Calculer un angle à l'aide de deux expressions du produit scalaire]
	Dans un repère orthonormé, on considère les points $A\coordl{-1}{2}$, $B\coordl{4}{1}$, $C\coordl{4}{3}$ et $D\coordl{2}{-1}$.
	\begin{center}
		\begin{tikzpicture}[x=9mm, y=9mm, >=Stealth, thick]
			\coordinate (A) at (-1,2);
			\coordinate (B) at (4,1);
			\coordinate (C) at (4,3);
			\coordinate (D) at (2,-1);
			\coordinate (O) at ({34/11},{13/11});
			\draw[xstep=1, ystep=1, gris!40, dashed, thin] (-1.2,-1.2) grid (5.5,4.2);
			\fill[bleu] (A) circle (2pt) node[left, bleu, fill=white, font=\footnotesize] {$A$};
			\fill[bleu] (B) circle (2pt) node[right, bleu, fill=white, font=\footnotesize] {$B$};
			\fill[rouge] (C) circle (2pt) node[above right, rouge, font=\footnotesize] {$C$};
			\fill[rouge] (D) circle (2pt) node[below left, rouge, font=\footnotesize, fill=white] {$D$};
			\fill[noir] (O) circle (2pt) node[above, xshift=-2, font=\footnotesize, fill=white] {$O$};
			\draw[->, gris] (-1.2,0) -- (5.5,0) node[above, font=\footnotesize] {$x$};
			\draw[->, gris] (0,-1.2) -- (0,4.2) node[right, font=\footnotesize] {$y$};
			\draw[->, noir, line width=1.1pt] (0,0) -- (1,0) node[midway, sloped, below, font=\footnotesize] {$\vec{\imath}$};
			\draw[->, noir, line width=1.1pt] (0,0) -- (0,1) node[midway, sloped, above, font=\footnotesize] {$\vec{\jmath}$};
			\draw[thin, noir] (0.25,0) -- (0.25,0.25) -- (0,0.25);
			\draw[dashed, gris] (A) -- (B);
			\draw[dashed, gris] (C) -- (D);
			\draw[->, bleu, line width=1.2pt] (A) -- (B) node[sloped, midway, above, bleu, font=\footnotesize] {$\vec{AB}$};
			\draw[->, rouge, line width=1.2pt] (C) -- (D) node[sloped, pos=0.2, above, rouge, font=\footnotesize] {$\vec{CD}$};
			\draw[->, vert, line width=1.2pt] (O) -- (B);
			\draw[thin, noir] ($(O) + (-14:0.55)$) arc[start angle=-14, end angle=-115, radius=0.55];
			\node[font=\footnotesize] at ($(O) + (-65:0.95)$) {$\widehat{BOD}$};
			\fill[bleu] (A) circle (2pt);
			\fill[bleu] (B) circle (2pt);
			\fill[rouge] (C) circle (2pt);
			\fill[rouge] (D) circle (2pt);
			\fill[noir] (O) circle (2pt);
		\end{tikzpicture}
	\end{center}
	Les droites $(AB)$ et $(CD)$ se coupent en $O$.

	Déterminer une valeur approchée au dixième de degré de l'angle $\widehat{BOD}$.

	\begin{correction}
		On vérifie d'abord que $O \in [AB] \cap [CD]$ : on trouve $O\coordl{\dfrac{34}{11}}{\dfrac{13}{11}}$, qui appartient aux deux segments.

		Les vecteurs $\vec{OB}$ et $\vec{AB}$ ont donc même sens, de même que $\vec{OD}$ et $\vec{CD}$ : l'angle $\widehat{BOD}$ est l'angle géométrique entre $\vec{AB}$ et $\vec{CD}$.

		On calcule $\vec{AB} \cdot \vec{CD}$ de deux façons :
		\begin{itemize}
			\item \textbf{Par les coordonnées :}

			      $\vec{AB}\coord{5}{-1}$ et $\vec{CD}\coord{-2}{-4}$, d'où :
			      \[
				      \vec{AB} \cdot \vec{CD} = 5 \times (-2) + (-1) \times (-4) = -6
			      \]

			\item \textbf{Par la formule du cosinus :}

			      $AB = \sqrt{26}$ et $CD = \sqrt{20}$, d'où :
			      \[
				      \begin{array}{rccl}
					      \vec{AB} \cdot \vec{CD} & = & \sqrt{26} \times \sqrt{20} \times \cos\pa{\widehat{BOD}} &     \\
					                              & = & \sqrt{520}\cos\pa{\widehat{BOD}}                         & =-6
				      \end{array}
			      \]
		\end{itemize}

		Donc :
		\[
			\cos\pa{\widehat{BOD}} = \dfrac{-6}{\sqrt{520}} = \dfrac{-3}{\sqrt{130}}
		\]

		À la calculatrice, $\widehat{BOD} \approx 105{,}3^\circ$ (angle obtus, cohérent avec le signe négatif du cosinus).
	\end{correction}
\end{methode}

